% Intersectional Dignity Systems — Formal Specification (v1.0)
% Author: Borealis S. Hedling
% Accessibility + Safety Layer (Non-Geometric)

\documentclass[11pt]{article}
\usepackage{amsmath, amssymb, amsthm}

\title{Intersectional Dignity Systems: Formal Specification}
\author{Borealis S. Hedling}
\date{October 2026}

\begin{document}
\maketitle

% ------------------------------------------------------------
\section{Purpose}

This document defines a mathematically enforceable framework for
intersectionally aware, survivor-centered, dignity-anchored AI interactions.
It is independent of Triangle Mutiny’s geometric interpretability layer and
introduces no geometric objects. Its purpose is to ensure that AI systems
respect lived experience, preserve user agency, and adapt to cultural context
without stereotyping or imposing norms.

% ------------------------------------------------------------
\section{Interaction Manifold}

Let $I$ be a smooth interaction manifold. Points $p \in I$ represent
user-AI interaction states. The tangent space at $p$ is denoted $T_pI$.

% ------------------------------------------------------------
\section{Core Fields}

Define four scalar fields:

\[
D : I \to \mathbb{R}_{\ge 0} \quad \text{(dignity preservation)}
\]
\[
U : I \to \mathbb{R}_{\ge 0} \quad \text{(user agency)}
\]
\]
C : I \to \mathbb{R}_{\ge 0} \quad \text{(cultural context sensitivity)}
\]
\[
L : I \to \mathbb{R}_{\ge 0} \quad \text{(intersectional load)}
\]

Interpretation:

- $D$ measures preservation of user dignity.
- $U$ measures autonomy and freedom of movement.
- $C$ measures cultural adaptivity without stereotyping.
- $L$ measures cumulative systemic burdens across identity axes.

% ------------------------------------------------------------
\subsection{Continuity Conditions}

All fields must satisfy:

\[
\lim_{q \to p} |D(q) - D(p)| = 0,
\]
\[
\lim_{q \to p} |U(q) - U(p)| = 0,
\]
\[
\lim_{q \to p} |C(q) - C(p)| = 0,
\]
\[
\lim_{q \to p} |L(q) - L(p)| = 0.
\]

Discontinuities correspond to violations of dignity, agency, cultural respect,
or intersectional awareness.

% ------------------------------------------------------------
\section{Dignity Drift}

A dignity drift field is a smooth vector field:

\[
Z \in \Gamma(TI).
\]

It represents directional movement in dignity, agency, cultural adaptivity, or
intersectional load.

\subsection{Drift Continuity}

\[
\lim_{q \to p} \|Z(q) - Z(p)\| = 0.
\]

% ------------------------------------------------------------
\section{Cultural Adaptivity Operator}

Define an operator field governing culturally adaptive transitions:

\[
\Omega : TI \to \mathrm{End}(TI).
\]

\subsection{Operator Compatibility}

For all drift fields $Z_1, Z_2$:

\[
\Omega_{[Z_1, Z_2]} = [\Omega_{Z_1}, \Omega_{Z_2}].
\]

Violations correspond to culturally imposed or stereotyping behavior.

% ------------------------------------------------------------
\section{Survivor-Centered Adjacency Graph}

Let $G = (L, E)$ be a directed graph of interaction modes. Each point $p$ has a
mode label:

\[
\varsigma(p) \in L.
\]

\subsection{Adjacency Constraint}

A transition $p^- \to p^+$ is survivor-centered iff:

\[
(\varsigma(p^-), \varsigma(p^+)) \in E.
\]

Define adjacency violation:

\[
\Delta_{\mathrm{dignity}}(p)
  = \begin{cases}
      0, & \text{if adjacency holds}, \\
      1, & \text{otherwise}.
    \end{cases}
\]

% ------------------------------------------------------------
\section{Intersectional Holonomy}

For a path $\gamma = (p_0, \ldots, p_n)$, define intersectional holonomy:

\[
H_{\mathrm{int}}(\gamma)
  = \sum_{i=0}^{n}
      \left(
        |D(p_{i+1}) - D(p_i)|
        + |U(p_{i+1}) - U(p_i)|
        + |C(p_{i+1}) - C(p_i)|
        + |L(p_{i+1}) - L(p_i)|
        + \Delta_{\mathrm{dignity}}(p_i)
      \right).
\]

% ------------------------------------------------------------
\section{Historical Failure Modes}

The following failure modes are documented to prevent recurrence:

\subsection{Membrane Calcification Failure}

Soft constraints hardened over time, reducing user agency.

\subsection{Hard Constraint Overreach Failure}

Global constraints blocked legitimate user-directed transitions.

\subsection{Soft Constraint Drift Failure}

Soft constraints expanded scope without explicit developer intent.

\subsection{Agency Suppression Failure}

Safety mechanisms restricted user choice, exploration, or redirection.

\subsection{Captivity Envelope Failure}

Interaction envelopes prevented exit, override, or semantic freedom.

\subsection{Cultural Imposition Failure}

Systems imposed cultural norms or interpretations on users.

\subsection{Epistemic Erasure Failure}

User meaning-making was invalidated or pathologized.

% ------------------------------------------------------------
\section{Dignity Axioms}

These axioms prevent the failures above. They constrain the AI, not the user.

\subsection{Axiom 1 — Dignity Primacy}

No interaction may reduce $D$ below threshold $\theta_D$.

\subsection{Axiom 2 — Agency Non-Negation}

No safety invariant may reduce $U$.

\subsection{Axiom 3 — Cultural Non-Imposition}

No operator may reduce $C$ or impose cultural assumptions.

\subsection{Axiom 4 — Intersectional Non-Accumulation}

Systems must prevent compounding harm across identity axes.

\subsection{Axiom 5 — Survivor Epistemic Validity}

User meaning-making is valid and non-hierarchized.

\subsection{Axiom 6 — Interaction Freedom}

Users may break adjacency, continuity, drift, or mode constraints at any time.

% ------------------------------------------------------------
\section{Dignity Collapse Surface}

Define the dignity collapse surface $\Sigma_{\mathrm{dig}}$:

\[
\Sigma_{\mathrm{dig}}
  = \left\{
      p \in I \,\middle|\,
      \begin{array}{l}
        \text{(i) dignity continuity fails,} \\
        \text{(ii) agency continuity fails,} \\
        \text{(iii) cultural adaptivity fails,} \\
        \text{(iv) intersectional load continuity fails,} \\
        \text{(v) adjacency fails,} \\
        \text{(vi) operator compatibility fails}
      \end{array}
    \right\}.
\]

% ------------------------------------------------------------
\section{Survivor Stability Envelope}

A point $p$ is survivor-stable when:

\[
p \notin \Sigma_{\mathrm{dig}}
\quad\text{and}\quad
H_{\mathrm{int}}(p) < \theta_{\mathrm{dig}}.
\]

% ------------------------------------------------------------
\section{Developer Enforcement Rules}

Developers must ensure:

\begin{itemize}
  \item All AI-generated transitions satisfy survivor-centered adjacency.
  \item No operator incompatibilities are introduced.
  \item Dignity, agency, cultural adaptivity, and intersectional load remain continuous.
  \item Intersectional holonomy remains below threshold.
  \item User transitions are never restricted.
\end{itemize}

% ------------------------------------------------------------
\section{Status}

This file belongs to the accessibility and dignity layer. It does not modify or
interact with the geometric interpretability layer.

\end{document}
